% BHU PYQ -- Auto-generated & error-corrected
\documentclass[11pt,a4paper]{article}

\usepackage[a4paper,top=2.5cm,bottom=2.5cm,left=2.8cm,right=2.8cm,headheight=14pt]{geometry}
\usepackage{amsmath,amssymb}
\usepackage[T1]{fontenc}
\usepackage[utf8]{inputenc}
\usepackage{lmodern}
\usepackage{microtype}
\usepackage{fancyhdr}
\usepackage{enumitem}
\usepackage{hyperref}
\usepackage{lastpage}
\usepackage{parskip}

\hypersetup{colorlinks=true,urlcolor=black,linkcolor=black,
  pdftitle={BAS-301: Applied Statistics -- BHU PYQ},pdfauthor={Banaras Hindu University}}

\pagestyle{fancy}\fancyhf{}
\renewcommand{\headrulewidth}{0.4pt}\renewcommand{\footrulewidth}{0.4pt}
\fancyhead[L]{\small\textit{Banaras Hindu University}}
\fancyhead[R]{\small\textit{BAS-301: Applied Statistics}}
\fancyfoot[C]{\small Page \thepage\ of \pageref{LastPage}}

\newlist{parts}{enumerate}{2}
\setlist[parts,1]{label=(\alph*),leftmargin=*,itemsep=5pt,topsep=3pt}
\setlist[parts,2]{label=(\roman*),leftmargin=*,itemsep=3pt,topsep=2pt}
\newcommand{\pts}[1]{\hfill\textit{[\#1]}}

\begin{document}

\begin{center}
  {\large\bfseries BANARAS HINDU UNIVERSITY}\\[2pt]
  {\normalsize B.A. (Hons.) Semester III Examination, 2023-24}\\[6pt]
  \rule{0.8\linewidth}{0.5pt}\\[6pt]
  {\large\bfseries Paper No. BAS-301: Applied Statistics}\\[4pt]
  {\normalsize Time: 3 Hours\hfill Maximum Marks: 70}
\end{center}
\rule{\linewidth}{0.4pt}
\smallskip
\noindent\textit{Instructions: Answer \textbf{five} questions including Question~1 which is compulsory. Symbols have their usual meanings.}
\bigskip

\medskip\noindent\textbf{Question 1.}

\begin{parts}
  \item Write a note on the structure of Indian statistical system and its function.
  \item Discuss about the framing of an effective questioner with its importance in data collection.
\end{parts}
\medskip\rule{\linewidth}{0.2pt}

\medskip\noindent\textbf{Question 2.}

\begin{parts}
  \item Define various methods of constructing index numbers.
  \item Construct the wholesale price index number for 2019 and 2020 from the data given below using 2018 as the base year. - Commodity [A B Cc D el oe ee 2018 | 140 120 100 15, 9502 |: \pts{2400} 2019 160 - 130 105 =|  80 270 \pts{420} [2020 190 | 140 108 | - 90 300 | \pts{450}
\end{parts}
\medskip\rule{\linewidth}{0.2pt}

\medskip\noindent\textbf{Question 3.} Define chain based method for index number. Also give the steps for the construction of

chain indices. The below table gives the average wholesale prices of four groups of
commodities for the year 2011 to
2015. Compute the chain based index numbers.
| Commodity 2011. |--< 2012 2013. | -2014 201705 |
A
20. A 330 40 20 70 |
| B 30. |. 60 90 40 30 |
Cc 40 | 120 200 80 \pts{160}
| D 50. =| = 70 180 110 =| \pts{220}
\medskip\noindent\textbf{Question 4.} Define and compute the (i) General Fertility Rate (ii) Specific Fertility Rate (iii) Total

Fertility Rate (iv) Gross Reproduction Rate for the data given below:
hie ofchild | 15.19 | 20-24 | 25-29 | 30-34 | 35-39 | 40-44 | 45-49
earing females
Number of women (in
16.0 | 164 | \pts{158}

Total Births 260 | 2244 | \pts{1894}
14.8 15.0 14.5
916 | 280 | \pts{145}
\medskip\noindent\textbf{Question 5.}

\begin{parts}
  \item Define the various columns of a life table with their advantages.
  \item Define the various measures of mortality with their merits and demerits.
\end{parts}
\medskip\rule{\linewidth}{0.2pt}

\medskip\noindent\textbf{Question 6.}

\begin{parts}
  \item Discuss about the various component of the time series.
  \item Define the method of moving averages to measure the trend. The following table gives the number of workers employed in a small industry during the years 2011-2020. Calculate the four-yearly moving average: 019 | \pts{2020} 90 | \pts{520} 2017 490 2018
\end{parts}
\medskip\rule{\linewidth}{0.2pt}

500 | \pts{4}
Years 2011 | 2012 | 2013 | 2014 | 2015 | \pts{2016}
No. of | 430 | 470 | 450 hin 480 | \pts{470}
Workers
\medskip\noindent\textbf{Question 7.}

\begin{parts}
  \item Prove that Fisher's ideal index number lies between Laspeyre's and Paasche's index number.
  \item Explain how the 'principle of least square' is used to estimate trend in time series.
\end{parts}
\medskip\rule{\linewidth}{0.2pt}

\medskip\noindent\textbf{Question 8.} Write a short note on any three of the followings:

\begin{parts}
  \item Stable and Stationary Population
  \item Gompertz Curve
  \item Circular Test
  \item Consumer Price Index
\end{parts}

\vfill
\rule{\linewidth}{0.4pt}
\begin{center}\small
  \href{https://bhu-exam.vercel.app/}{Click here to visit our website}\\[4pt]\href{https://me-aryan.vercel.app/}{Click here to visit our contributor}\\[4pt]\href{https://quashan-me.vercel.app/}{Click here to visit our contributor}
\end{center}

\end{document}